% Source: https://artofproblemsolving.com/community/c3h1379390_triangular_prism_surface_area
% Stable identity: aops-topic-1379390-q01
The surface of a prism consists of the two triangular bases and three rectangular lateral faces, one for each side of the triangle, each as long as the prism.

\begin{asy}
usepackage("amsmath");
size(6cm);
defaultpen(fontsize(10pt));
path ra(pair P,pair u,pair v,real s=0.2){return P+s*unit(u)--P+s*unit(u)+s*unit(v)--P+s*unit(v);}
void tk(pair P,pair Q,int n,real t=0.5,real h=0.08,real g=0.07,pen p=currentpen){pair m=P+t*(Q-P),u=unit(Q-P),v=rotate(90)*u;for(int i=0;i<n;++i){pair c=m+(i-(n-1)/2)*g*u;draw(c-h*v--c+h*v,p);}}
pair v=(4.33,2.5);
pair P0=(0,0),P1=(4,0),P2=(0,3);
draw(P0+v--P1+v,dashed); draw(P0+v--P2+v,dashed); draw(P0--P0+v,dashed);
draw(P0--P1--P2--cycle,linewidth(0.9));
draw(P1--P1+v--P2+v--P2,linewidth(0.9));
fill(P0--P1--P2--cycle,lightgray+opacity(0.5));
draw(ra(P0,P1-P0,P2-P0,0.35));
label("$4$",(P0+P1)/2,S); label("$3$",(P0+P2)/2,W); label("$5$",P2+0.3*(P1-P2),NE);
label("$10$",P1+v/2,SE);
label("\mbox{}",(min(currentpicture,user=true)+max(currentpicture,user=true))/2,black); // keeps the text after the figure black when LaTeX typesets the labels inline
\end{asy}

\begin{itemize}
\item The hypotenuse of the base is $\sqrt{3^2+4^2}=5$ cm.
\item Each base has area $\dfrac{1}{2}\cdot3\cdot4=6$ cm$^2$, so the two bases give $12$ cm$^2$.
\item The lateral faces are $3\times10$, $4\times10$ and $5\times10$ rectangles, with total area $10(3+4+5)=120$ cm$^2$, that is, the perimeter of the base times the length.
\end{itemize}
The total surface area is $12+120=\boxed{132\text{ cm}^2}$.

In general, for a right triangular prism of length $\ell$ whose base has area $B$ and perimeter $P$, the surface area is $2B+P\ell$.
